% Essential packages
\usepackage{graphicx}
\usepackage{comment}
\usepackage{amssymb}
\usepackage{amsmath}
\usepackage{amsthm}
\usepackage{graphicx}
\usepackage{float}
\usepackage{comment}
\usepackage{caption}
\usepackage{multirow}
\usepackage{listings}
\usepackage{xcolor}
\usepackage[hidelinks]{hyperref}
\usepackage{enumitem}
\usepackage{titlesec}
\usepackage{algorithm}
\usepackage{algpseudocode}
\usepackage{tcolorbox}
\usepackage{fancyhdr}
\usepackage{mdframed}
\usepackage{tikz}
\usepackage{pgfplots}
\usepackage{parskip}

% Page margins: slightly wider left/right than the article default
\usepackage{geometry}
\geometry{hmargin=2in, vmargin=1in}

% Command to get the current month name
\newcommand{\MonthName}{
  \ifcase\month\or January\or February\or March\or April\or May\or June\or July\or August\or September\or October\or November\or December\fi}

% Configurazione di codeblocks con il pacchetto listings
\lstset{
  basicstyle=\ttfamily\small,
  keywordstyle=\color{blue},
  commentstyle=\color{green!50!black},
  stringstyle=\color{red},
  showstringspaces=false,
  breaklines=true,
  frame=single, % Adds a border around the code
  numbers=left, % Line numbers on the left
  numberstyle=\tiny\color{gray}, % Style of line numbers
  stepnumber=1, % Number of lines to skip between numbers
  numbersep=-5pt, % Distance between line numbers and code
  tabsize=2, % Tab size
  captionpos=b, % Position of the caption
}

% Header configuration
\pagestyle{fancy}
\fancyhf{} % Clean header and footer
\fancyhead[L]{\leftmark} % Current section on the top left
\fancyhead[R]{\thepage} % Page number on the right

% Defines a command for framed definitions with a blue background and a title
\newcommand{\definition}[2]{%
\begin{mdframed}[
    hidealllines=false,
    linecolor=blue!80!black,
    linewidth=1pt,
    backgroundcolor=blue!5,
    roundcorner=5pt,
    frametitle={#1},
    frametitlebackgroundcolor=blue!80!black,
    frametitlefont=\color{white}\bfseries,
    innerbottommargin=10pt,
    skipabove=10pt,
    skipbelow=10pt
  ]
  {#2}
  \end{mdframed}
}

% Defines a command for framed Note boxes with green background and a title
\newcommand{\note}[1]{%
\begin{mdframed}[
    hidealllines=false,
    linecolor=green!80!black,
    linewidth=1pt,
    backgroundcolor=green!5,
    roundcorner=5pt,
    frametitle={Note},
    frametitlebackgroundcolor=green!80!black,
    frametitlefont=\color{white}\bfseries,
    innerbottommargin=10pt,
    skipabove=10pt,
    skipbelow=10pt
  ]
  {#1}
  \end{mdframed}
}

% Make proof keyword bold
\makeatletter
\def\proofname{\bfseries Proof}
\makeatother
% PGFPlots settings
\pgfplotsset{compat=1.18}
